\input{../header}
\hyphenpenalty=5000

\rhead{Your name: \rule{8cm}{0.15mm}}

\begin{document}

\everymath{\displaystyle}

\section{Learning target A4, version 1}

Use definite integrals to find the mass and center of mass of an object.

\hrulefill

Say you have a thin rod that is 6 cm long, and its left end lies at $x = 0$ on the horizontal axis.  Suppose that the rod's density is distributed according to the function $\rho(x) = 2e^{-(x-4)^2}$, measured in grams per cm.

\begin{enumerate}
	\item From the graph of $\rho(x)$ (without doing any computations), do you expect the rod's center of mass, $\overline{x}$, to lie at $x = 3$, left of $x = 3$, or right of $x = 3$?  Why?
	
	\vfill
	
	\item Find the total mass $M$ and center of mass $\overline{x}$ by writing down appropriate definite integrals and evaluating them using technology.  Clearly identify your results.

	\vfill
	\vfill

	\item How would the rod's mass and center of mass change (if at all) if the density function was instead $\rho(x) = 2e^{-(x-2)^2}$? Don't do any computations, but write a sentence or two to discuss.

	\vfill
	
\end{enumerate}

\end{document}